% !PRINTED 81
\weberpage{81}
%<<CONT
then the relation between the $\alpha$ and the $\beta$ is expressed by the
equation

\weq{1}{\omega_1^{\alpha_1\beta_1}\,\omega_2^{\alpha_2\beta_2}\,\dots\,\omega_\mu^{\alpha_\mu\beta_\mu}=1}

The invariants of the group $\mathfrak{B}$, which by \S.\,10 are all prime powers, we shall denote by
\[
i_1, i_2, \dots i_\nu \quad (\text{invariants of } \mathfrak{B})
\]
We take these invariants to be arbitrarily given
prime powers, and denote by $J$ their least common multiple, so that
\weq{2}{J=i_1i_1'=i_2i_2'=\dots=i_\nu i_\nu'}
can be set.

Moreover, $\mathfrak{A}$ shall be assumed to be a primary divisor of $\mathfrak{N}$, which by \S.\,19, 8. is equivalent to the condition:

\gesperrt{IV. None of the indices $\beta$ shall, for all numbers $b$, be divisible by $q$ when the corresponding prime $q$ divides $m$ more than once, or by $q-1$ when $q$ divides $m$ only once.}

By the definition of the invariants, the group $\mathfrak{B}$ must have a basis whose elements
\[
\begin{aligned}
&\phantom{\text{of the degrees}} \quad g_1, g_2, \dots g_\nu \quad (\text{basis of } \mathfrak{B})\\
&\text{of the degrees} \quad i_1, i_2, \dots i_\nu
\end{aligned}
\]
are, so that every number $b$ is contained once and only once in the form
\weq{3}{b\equiv g_1^{x_1}g_2^{x_2}\dots g_\nu^{x_\nu}\ (\mathrm{mod}\ m)}
when the exponents $x_1, x_2, \dots x_\nu$ each run through a complete system of residues with respect to the moduli $i_1, i_2, \dots i_\nu$.

All quantities of the form (3), if $g_1, g_2, \dots g_\nu$ are arbitrary numbers, certainly form a group under unrestricted variation of the integer exponents $x$. But if the $g_h$ are to be the elements of a basis of this group, and $i_h$ their degrees, then $1\equiv g_1^{x_1}g_2^{x_2}\dots g_\nu^{x_\nu}\ (\mathrm{mod}\ m)$ may be satisfied only when $x_h$ is divisible by $i_h$; hence:

\gesperrt{1. The necessary and sufficient condition that $g_1, g_2, \dots g_\nu$ be the elements of a basis of a group $\mathfrak{B}$ of the degrees $i_1, i_2, \dots i_\nu$ is that the congruence}
\weq{4}{g_1^{x_1}g_2^{x_2}\dots g_\nu^{x_\nu}\equiv 1\ (\mathrm{mod}\ m)}
%CONT>>
